\documentclass{article}

\usepackage{amsmath,amssymb}
\usepackage{xurl}
\usepackage[hidelinks]{hyperref}
\setlength{\emergencystretch}{2em}

% Shared commands for notes in docs/paper-gaps/.

\DeclareUnicodeCharacter{2080}{\ifmmode{}_{0}\else\textsubscript{0}\fi}
\DeclareUnicodeCharacter{2081}{\ifmmode{}_{1}\else\textsubscript{1}\fi}
\DeclareUnicodeCharacter{2082}{\ifmmode{}_{2}\else\textsubscript{2}\fi}
\DeclareUnicodeCharacter{2099}{\ifmmode{}_{n}\else\textsubscript{n}\fi}

\newcommand{\Lean}{\textsc{Lean}}
\ExplSyntaxOn
\NewDocumentCommand{\leanid}{m}{
  \ifmmode
    \textsf{\detokenize{#1}}
  \else
    \group_begin:
    \urlstyle{sf}
    \tl_set:Nn \l_tmpa_tl {#1}
    \tl_replace_all:Nnn \l_tmpa_tl {\_} {_}
    \exp_args:NV \path \l_tmpa_tl
    \group_end:
  \fi
}
\ExplSyntaxOff
\providecommand{\tr}{\operatorname{tr}}
\newcommand{\tnleanIssueBase}{https://github.com/LionSR/TNLean/issues/}
\newcommand{\tnleanPrBase}{https://github.com/LionSR/TNLean/pull/}
\newcommand{\tnleanDocBase}{https://sirui-lu.com/TNLean/docs/}
\newcommand{\ghissue}[1]{\href{\tnleanIssueBase#1}{GitHub issue \##1}}
\newcommand{\ghpr}[1]{\href{\tnleanPrBase#1}{PR \##1}}
\newcommand{\doclink}[1]{\href{\tnleanDocBase#1.html}{\path{#1}}}

% Dirac notation and the identity operator, as in blueprint/src/macros/common.tex.
% \providecommand keeps note-local definitions authoritative.
\usepackage{dsfont}
\providecommand{\ket}[1]{|#1\rangle}
\providecommand{\bra}[1]{\langle #1|}
\providecommand{\braket}[2]{\langle #1 | #2 \rangle}
\providecommand{\ketbra}[2]{|#1\rangle\!\langle #2|}
\providecommand{\Id}{\mathds{1}}
\providecommand{\one}{\mathds{1}}
\providecommand{\C}{\mathbb{C}}
\providecommand{\MN}[1]{M_{#1}(\C)}
\providecommand{\MD}{M_D(\C)}

% Verdict marker: \gapnote{<kind>}{<status>}. Kind is the policy
% classification (clarification, local-correction, scope-restriction,
% unfaithful, false-source, open-gap); status is open, wip (elimination
% underway), resolved, or historical. Machine-read by the site index;
% prints a verdict line.
\newcommand{\gapnote}[2]{\par\noindent\textbf{Verdict:} #1 (#2).\par}


% Fallbacks keep this note compilable when a different command.tex is found first.
\providecommand{\leanid}[1]{\textsf{\detokenize{#1}}}
\providecommand{\gapnote}[2]{\par\noindent\textbf{Verdict:} #1 (#2).\par}

\title{Lattice Points and the Stopping Scale in Belt and Replacement Templates}
\author{}
\date{2026-10-16}

\begin{document}
\maketitle
\gapnote{scope-restriction}{resolved}

\section{The source passage}

Section~11 of the September~24, 2026 two-dimensional area-law
manuscript~\cite{OpenAI2026AreaLaw} (\path{10-geometry.tex}, lines~731--740)
assigns to a belt tile of side scale $u=t_k$ the parameters $s_0=u$ and
$n=\lceil C_Bu\rceil$, and to a replacement with removal half-side $u$ the
parameters $s_0=2u$ and $n=\lceil C_Bu\rceil$. Each region is a union of at
most eight closed triangles, and $C_B$ is chosen large enough for every such
template to satisfy $n\ge C_{\rm tpl}N_0(s_0+1)$ and
$\operatorname{diam}_\infty T\le n$.

\section{The formal statements}

The formal results\footnote{\leanid{TNLean.PEPS.AreaLaw.Geometry.initialBelt_exists_template},
\leanid{TNLean.PEPS.AreaLaw.Geometry.repairReplacement_exists_template}.}
carry two premises that the passage does not print.
\begin{enumerate}
  \item The region contains a lattice point. Definition~9.3
        (\path{08-scanner.tex}, lines~546--564) requires a template to be a
        nonempty set, so a region with no lattice point has no template. The
        source uses these templates only for identifiers whose final regions
        meet the cut $A$ (\path{10-geometry.tex}, lines~776--779). Such a
        region contains a lattice point.
  \item For replacements, the stopping scale satisfies $S_*\ge1$. Then every
        nonterminal removal half-side is a natural power of two, so $s_0=2u$
        is a positive integer. The source chooses $S_*$ sufficiently large
        (lines~380--394), and requires a terminal half-side produced from
        $S\ge S_*$ to be at least two (lines~549--552). That choice implies
        $S_*\ge1$.
\end{enumerate}

\section{Elimination}

Neither premise restricts the use of these results in the proof of the
two-family proposition. The final assembly applies the theorems only to
retained regions, which meet $A$, and with a stopping scale chosen at least
one. It discharges both premises there, and no unrestricted statement is
needed. The unrestricted reading of the first premise
is vacuous, because Definition~9.3 has no empty template.

The first premise is discharged by the lemma on lattice points in retained
births.\footnote{\leanid{TNLean.PEPS.AreaLaw.Geometry.repairStage_exists_integerPoint_mem_birth_of_oldInclusion},
\leanid{TNLean.PEPS.AreaLaw.Geometry.repairStage_retained_exists_integerPoint_mem_birth}.}
If an identifier keeps a current region meeting $A$ at a later generation,
then its birth, which retention leaves unchanged, contains a lattice point:
a point of $A$ is a lattice point, and every current region lies in its
birth. The second premise is absorbed into the stopping-scale threshold by
taking the threshold at least one, in the theorem that supplies templates for
all actual births.\footnote{\leanid{TNLean.PEPS.AreaLaw.Geometry.exists_uniform_repairStage_birth_templates}.}

\section{Resolution}

Both premises are discharged in the proof of the two-family
proposition.\footnote{\leanid{TNLean.PEPS.AreaLaw.twoFamilyDecomposition}.}
Each retained piece there has a current region containing a point of $A$; since
current regions lie in their births, that point is a lattice point of the birth,
which gives the first premise. The stopping scale is chosen at least the
template threshold, and that threshold is at least one, which gives the second. The
templates of the belt tiles and replacements are therefore used only under
hypotheses the source construction provides.

\bibliographystyle{alpha}
\bibliography{references}

\end{document}
